\documentclass[11pt,a4paper]{article}
\usepackage[T1]{fontenc}
\usepackage{microtype}
\usepackage{amsmath,amssymb,amsthm,mathtools}
\usepackage{mathrsfs}
\usepackage{aliascnt}
\usepackage{booktabs,array,longtable}
\usepackage{geometry}
\usepackage{xcolor}
\usepackage{hyperref}
\usepackage[nameinlink,noabbrev]{cleveref}

\geometry{top=2.45cm,bottom=2.45cm,left=2.75cm,right=2.75cm}
\setlength{\emergencystretch}{2em}
\raggedbottom
\hypersetup{
  colorlinks=true,
  linkcolor=blue!55!black,
  citecolor=blue!55!black,
  urlcolor=blue!55!black,
  pdfauthor={Akari Hayami}
}

\theoremstyle{plain}
\newtheorem{theorem}{Theorem}[section]
\newaliascnt{proposition}{theorem}
\newtheorem{proposition}[proposition]{Proposition}
\aliascntresetthe{proposition}
\newaliascnt{lemma}{theorem}
\newtheorem{lemma}[lemma]{Lemma}
\aliascntresetthe{lemma}
\newaliascnt{corollary}{theorem}
\newtheorem{corollary}[corollary]{Corollary}
\aliascntresetthe{corollary}
\theoremstyle{definition}
\newaliascnt{definition}{theorem}
\newtheorem{definition}[definition]{Definition}
\aliascntresetthe{definition}
\newaliascnt{example}{theorem}
\newtheorem{example}[example]{Example}
\aliascntresetthe{example}
\theoremstyle{remark}
\newaliascnt{remark}{theorem}
\newtheorem{remark}[remark]{Remark}
\aliascntresetthe{remark}
\crefname{proposition}{proposition}{propositions}
\Crefname{proposition}{Proposition}{Propositions}
\crefname{lemma}{lemma}{lemmas}
\Crefname{lemma}{Lemma}{Lemmas}
\crefname{corollary}{corollary}{corollaries}
\Crefname{corollary}{Corollary}{Corollaries}
\crefname{definition}{definition}{definitions}
\Crefname{definition}{Definition}{Definitions}
\crefname{example}{example}{examples}
\Crefname{example}{Example}{Examples}
\crefname{remark}{remark}{remarks}
\Crefname{remark}{Remark}{Remarks}
\numberwithin{equation}{section}

\newcommand{\N}{\mathbb N}
\newcommand{\R}{\mathbb R}
\newcommand{\Z}{\mathbb Z}
\newcommand{\C}{\mathbb C}
\newcommand{\T}{\mathbb T}
\newcommand{\F}{\mathbb F}
\newcommand{\kk}{\Bbbk}
\newcommand{\norm}[1]{\left\lVert#1\right\rVert}
\newcommand{\Aex}{\mathcal A_{\mathrm{ex}}}
\newcommand{\Hex}{H_{\mathrm{ex}}}
\newcommand{\Nex}{\mathcal N_{\mathrm{ex}}}
\DeclareMathOperator{\im}{im}
\DeclareMathOperator{\id}{id}
\DeclareMathOperator{\dist}{dist}
\DeclareMathOperator{\reach}{reach}
\DeclareMathOperator{\Fill}{Fill}
\DeclareMathOperator{\Fillw}{\widehat{Fill}}
\DeclareMathOperator{\End}{End}
\DeclareMathOperator{\supp}{supp}
\DeclareMathOperator{\rank}{rank}
\DeclareMathOperator{\ob}{ob}
\usepackage{graphicx}
\usepackage{tikz}
\usepackage{pgfplots}
\usepackage{caption}
\usepackage[section]{placeins}
\usetikzlibrary{arrows.meta,positioning,calc,fit,backgrounds}
\pgfplotsset{compat=1.18}
\definecolor{viewblue}{HTML}{245A81}
\definecolor{viewteal}{HTML}{16766D}
\definecolor{vieworange}{HTML}{AE591E}
\definecolor{viewgray}{HTML}{53616C}
\tikzset{
  >=Latex,
  every picture/.style={font=\small},
  vbox/.style={draw=viewblue,rounded corners=2pt,fill=viewblue!5,
    align=center,inner sep=7pt,line width=.7pt},
  vinput/.style={vbox,draw=vieworange,fill=vieworange!6},
  vresult/.style={vbox,draw=viewteal,fill=viewteal!6},
  varrow/.style={->,line width=.8pt,draw=viewblue},
  vconditional/.style={->,dashed,line width=.8pt,draw=vieworange},
  vlabel/.style={font=\footnotesize,align=center,fill=white,inner sep=2pt},
  vpoint/.style={circle,fill=viewblue,inner sep=2.5pt},
  vseed/.style={circle,draw=vieworange,fill=vieworange!20,inner sep=3pt}
}
\pgfplotsset{viewaxis/.style={
  axis lines=left,axis line style={viewgray},
  tick label style={font=\scriptsize},label style={font=\small},
  legend style={font=\scriptsize,draw=none,fill=white},
  grid=major,grid style={viewgray!12},
  every axis plot/.append style={line width=.85pt},
  scaled ticks=false
}}
\captionsetup[figure]{font=small,labelfont=bf,labelsep=period}
\renewcommand{\topfraction}{.9}
\renewcommand{\bottomfraction}{.85}
\renewcommand{\textfraction}{.08}
\renewcommand{\floatpagefraction}{.75}
\setcounter{topnumber}{3}
\setcounter{bottomnumber}{2}
\setcounter{totalnumber}{4}
\newcommand{\ViewHeading}[1]{\par\smallskip{\small\bfseries #1}\par\smallskip}
\hypersetup{pdftitle={Foundations of Information Exclusion and Relative Obstruction}}

\title{\textbf{Foundations of Information Exclusion}\\
\textbf{and Relative Obstruction}\\
\large Finite Labels, Persistent Filling, and Local Extension}
\author{Akari Hayami}
\date{Reorganized foundations manuscript --- 16 September 2026}

\begin{document}
\maketitle

\begin{abstract}
We give a limited algebraic meaning to information exclusion.  In a finite
chain complex over $\F_Q$, a boundary-writing operation produces exactly
the boundary subspace, and the quotient labels form homology.  A uniformly
chosen homology label has $\beta_n\log_2 Q$ bits of entropy.  This counts
all quotient labels, including zero; it is not a channel-capacity theorem.
For a finite inclusion filtration, the same boundary equation defines a
precise filling threshold.  A three-edge, one-face example exhibits the
death of one nonzero label.  Separately, a local section fails to extend
exactly when its relative connecting class is nonzero.

These are standard algebraic mechanisms assembled with their hypotheses
kept explicit.  We distinguish support costs from normed costs, fixed
boundary spaces from parameterized operators, and qualitative obstruction
tests from metric recurrence.  The companion \emph{Arithmetic Local-System
Filling Spectra} constructs one finite-at-each-scale quantitative model
linking reciprocal filling cost to Diophantine recurrence and realizing
its parameter spectrum on positive-reach carriers.  That concrete result
does not supply a canonical bridge for arbitrary filtered complexes.
We state a conditional dimension-and-nonextension interface with explicit
metric, coefficient, and seed inputs.  A geometric appendix proves the
standard sub-reach retraction for general closed positive-reach sets,
without asserting an automatic finite cellular model.
\end{abstract}

\tableofcontents
\clearpage

\section{What the three manuscripts establish}
\label{sec:scope}

The purpose of this foundations paper is to specify the objects before
transferring conclusions between them.  A homology class, a relative
extension obstruction, and a large norm required to solve a boundary
equation are not interchangeable descriptions of one event.

The series is organized by its logical requirements rather than the
historical order in which the examples were developed:
\begin{enumerate}
 \item This paper establishes finite quotient-label counts, persistent
 boundary filling, and the relative extension criterion.
 \item \emph{Arithmetic Local-System Filling Spectra} supplies an
 explicitly marked, normed arithmetic model and a geometric realization
 of its selected parameter locus.
 \item \emph{Minimal Marked Response Modules for Relative Obstruction
 Data} closes a supplied finite seed under supplied endomorphisms.  It is
 an algebraic companion, not a further automatic stage of the arithmetic
 construction.
\end{enumerate}

The concrete bridge in the second paper is visible from the outset.  For
$x\in\T^d$ it uses the row
\[
 v_q(x)=(e^{2\pi\mathrm i qx_j}-1)_{j=1}^d
\]
and obtains the exact minimum Euclidean filling cost
$1/\norm{v_q(x)}_2$, interpreted as $+\infty$ at a zero row.  The
recurrence selection $\norm{v_q(x)}_2<Cq^{-\alpha}$ infinitely often
has dimension $(d+1)/(1+\alpha)$ for $\alpha>1/d$.  The coefficient
convention, arithmetic maps, metric comparisons, and classical dimension
input are all supplied in that paper.  None is extracted from the
finite-field model below.

We use $Q$ for the order of a finite field and reserve $q$ for arithmetic
scale in the companion, avoiding two meanings of the same parameter.

\Cref{fig:series-overview} separates the supplied inputs of the three
papers; the cross-series dictionary in \cref{fig:series-dictionary}
also distinguishes their outputs and indexing directions.
\begin{figure}[tbp]
\centering
\begin{tikzpicture}
\node[vbox,text width=12.4cm] (found) at (0,0)
  {\textbf{Foundations: qualitative algebra with specified data}};
\node[vbox,text width=5.3cm] (labels) at (-3.55,-1.55)
  {Finite-field chains and inclusions\\
   $Z_n/B_n$, label counts, filling thresholds};
\node[vbox,text width=5.3cm] (relative) at (3.55,-1.55)
  {Relative cochains and local sections\\
   $\delta(s)$, quotient and seed};
\draw[varrow] (found.south -| labels.north) -- (labels);
\draw[varrow] (found.south -| relative.north) -- (relative);
\node[vinput,text width=5.3cm] (arinput) at (-3.55,-3.3)
  {\textbf{Additional arithmetic model}\\
   Characters, norms, power maps, metric, carrier};
\node[vinput,text width=5.3cm] (marks) at (3.55,-3.3)
  {\textbf{Additional response marks}\\
   Indexed $G_j$ with $G_j(\mathsf B)\subseteq\mathsf B$};
\node[vresult,text width=5.3cm] (arithmetic) at (-3.55,-5.15)
  {\textbf{Arithmetic Filling Spectra}\\
   Quantitative recurrence and geometric transport};
\node[vresult,text width=5.3cm] (response) at (3.55,-5.15)
  {\textbf{Minimal Response Modules}\\
   $\Nex=\mathscr R_t\,\operatorname{im}\ob$};
\draw[varrow] (arinput) -- (arithmetic);
\draw[varrow] (marks) -- (response);
\draw[varrow] (relative.east) -- ++(.65,0) |- (response.east)
  node[pos=.5,vlabel,rotate=90] {quotient and seed};
\node[align=center,font=\footnotesize,text width=12.5cm] at (0,-6.4)
  {No automatic arrow from finite labels to arithmetic data,\\
   or from arithmetic scale maps to fixed-quotient endomorphisms.};
\end{tikzpicture}
\caption{Logical inputs of the three manuscripts. Solid arrows indicate
supplied constructions, not a universal pipeline. The arithmetic example
and the response marks are additional inputs; neither is selected by
relative exactness.}
\label{fig:series-overview}
\end{figure}

\section{Boundary writing and finite quotient labels}
\label{sec:labels}

Fix a finite field $\F_Q$ and a finite-dimensional chain complex
\[
 \cdots\longrightarrow C_{n+1}\xrightarrow{\partial_{n+1}}C_n
           \xrightarrow{\partial_n}C_{n-1}\longrightarrow\cdots,
 \qquad\partial_n\partial_{n+1}=0.
\]
Set
\[
 Z_n=\ker\partial_n,\qquad B_n=\im\partial_{n+1},\qquad
 H_n=Z_n/B_n,\qquad\beta_n=\dim_{\F_Q}H_n.
\]

\begin{definition}[Boundary-write model]\label{def:writing}
An $n$-cycle is \emph{boundary-writable} when it belongs to $B_n$.
Two cycles encode the same \emph{quotient label} when their difference
belongs to $B_n$.  A nonzero quotient label is an \emph{exclusion label}
for this specified writing operation.
\end{definition}

This defines a finite algebraic recording convention.  It does not assert
that every physical writing process is a boundary map, or that a nonzero
homology class is information that cannot be represented by any other
operation.

\begin{theorem}[Exact finite-label count]\label{thm:labels}
There are $Q^{\beta_n}$ quotient labels and $Q^{\beta_n}-1$ nonzero
exclusion labels.  If $L_n$ is uniform on all of $H_n$, then
\[
 \mathsf H(L_n)=\log_2|H_n|=\beta_n\log_2 Q.
\]
When $\beta_n>0$, a variable uniform only on the nonzero labels instead
has entropy $\log_2(Q^{\beta_n}-1)$.
\end{theorem}

\begin{proof}
A $\beta_n$-dimensional vector space over $\F_Q$ has
$Q^{\beta_n}$ elements, exactly one of which is zero.  The entropy of a
uniform variable on a finite set of size $m$ is $\log_2m$.
\end{proof}

For $\beta_n=0$ the nonzero-label set is empty, so there is no uniform
probability distribution on it.  The uniform distribution on $H_n$ is
still defined and has entropy zero.  Thus the label count and the choice
of probability space must both be stated.

\begin{remark}[Entropy is not capacity]
The displayed entropy concerns a specified uniform experiment on a
quotient set.  A channel capacity would additionally require a channel
law and a class of admissible input distributions.  No operational coding
rate, thermodynamic entropy, or erasure energy is asserted here.
\end{remark}

\subsection{Support filling cost}

The quotient and its sampling convention are illustrated in
\cref{fig:foundations-labels}.
\begin{figure}[tbp]
\centering
\begin{minipage}[t]{.46\linewidth}
\centering\ViewHeading{(a) Subspaces and quotient classes}
\begin{tikzpicture}[scale=.85]
\draw[rounded corners,draw=viewgray,fill=viewgray!3] (-3,-1.8) rectangle (3,1.8);
\node at (2.5,1.45) {$C_n$};
\draw[draw=viewblue,fill=viewblue!6] (0,0) ellipse (2.35 and 1.25);
\node at (1.85,.75) {$Z_n$};
\draw[draw=viewteal,fill=viewteal!15] (-.9,0) ellipse (.7 and .9);
\node at (-.9,0) {$B_n$};
\draw[draw=vieworange,dashed,rounded corners] (.05,-.85) rectangle (1.25,.65);
\node[vpoint,label=above:$z$] at (.4,.3) {};
\node[vpoint,label=below:{$z+b$}] at (.8,-.45) {};
\node[font=\footnotesize] at (0,-2.25) {$b\in B_n\quad\Longrightarrow\quad[z+b]=[z]$};
\end{tikzpicture}
\end{minipage}\hfill
\begin{minipage}[t]{.50\linewidth}
\centering\ViewHeading{(b) An illustrative $H_n\cong\F_2^2$}
\begin{tikzpicture}
\node[vseed,label=below:$0$] at (0,0) {};
\node[vpoint,label=below:$a$] at (1.4,0) {};
\node[vpoint,label=below:$b$] at (2.8,0) {};
\node[vpoint,label=below:{$a+b$}] at (4.2,0) {};
\node[vbox,text width=5.6cm] at (2.1,-1.25)
  {All four labels, uniform: $2$ bits};
\node[vinput,text width=5.6cm] at (2.1,-2.3)
  {Three nonzero labels, uniform: $\log_2 3$ bits};
\end{tikzpicture}
\end{minipage}
\caption{The quotient identifies cycles differing by a boundary. The
regions in (a) depict algebraic containment, not a topology on a finite
field. Panel (b) is an illustrative two-dimensional quotient, separate
from the one-dimensional triangle example. Entropy depends on the stated
sample space and uniform distribution; it is not channel capacity.}
\label{fig:foundations-labels}
\end{figure}

Choose a basis of $C_{n+1}$.  For $z\in Z_n$, define
\begin{equation}\label{eq:support-cost}
 \Fill^{\mathrm{supp}}(z)
 =\min\{|\supp c|:c\in C_{n+1},\ \partial_{n+1}c=z\},
 \qquad\min\varnothing=+\infty.
\end{equation}

\begin{proposition}[Exact boundary obstruction]\label{prop:boundary}
For every $z\in Z_n$,
\[
 \Fill^{\mathrm{supp}}(z)=+\infty
 \iff z\notin B_n
 \iff[z]\ne0\text{ in }H_n.
\]
\end{proposition}

\begin{proof}
The feasible set is nonempty precisely when $z$ lies in the image of
$\partial_{n+1}$.  A cycle has zero quotient class precisely when it lies
in that image.  When feasible, the finite chain space supplies a minimum.
\end{proof}

The finite value of the cost is basis dependent.  For example, replacing
the coordinate basis of $\F_2^2$ by $((1,1),(0,1))$ changes the support
size of $(1,1)$ from two to one.  The finite/infinite distinction is
basis independent because it only asks whether the boundary equation has
a solution.  The Euclidean cost of the arithmetic companion is a
different functional over a different coefficient field, not an
unstated continuation of \cref{eq:support-cost}.

\section{Finite filtrations and persistent filling thresholds}
\label{sec:persistence}

Let $K_{t_0}\subseteq\cdots\subseteq K_{t_m}$ be a finite filtration of
finite simplicial complexes over $\F_Q$.  Use compatible oriented
simplicial bases.  Choose $z\in Z_n(K_{t_i})$ and carry that same chain
through subsequent inclusions.  The chosen starting level $t_i$ need not
be the birth of a barcode interval or the earliest level containing a
representative of its class.

\begin{definition}[Filling threshold]\label{def:threshold}
The threshold of this carried cycle is
\[
 d(z)=\min\{t_j:j\geq i,\ z\in B_n(K_{t_j})\},
\]
with $d(z)=+\infty$ when the set is empty.  Its exclusion interval is the
set of available levels $t_j\geq t_i$ satisfying $t_j<d(z)$.
\end{definition}

\begin{proposition}[Threshold and persistent labels]
\label{prop:persistence}
For every $j\geq i$,
\[
 \Fill^{\mathrm{supp}}_{t_j}(z)=+\infty\iff t_j<d(z).
\]
Once finite, the support cost cannot increase along the filtration.
Furthermore, a uniform variable on
\[
 P_n^{i,j}=\im\bigl(H_n(K_{t_i};\F_Q)\to H_n(K_{t_j};\F_Q)\bigr)
\]
has entropy $\dim_{\F_Q}P_n^{i,j}\cdot\log_2Q$.
\end{proposition}

\begin{proof}
A filler at one level remains a filler at every later level because the
chain maps are inclusions.  Thus fillability is monotone, and
\cref{prop:boundary} gives the threshold equivalence.  Its old support
size remains unchanged in the larger compatible basis, so the minimum
cost cannot increase.  Apply \cref{thm:labels} to the vector space
$P_n^{i,j}$ to obtain the entropy statement.
\end{proof}

This uses the familiar finite persistence setting
\cite{zomorodian-carlsson}, but needs neither an interval-decomposition
theorem nor a stability or sampling theorem.  It does not extend the same
monotonicity to arithmetic power maps or to arbitrarily rescaled norms.

\begin{example}[One label killed by one face]\label{ex:triangle}
Over $\F_2$, let $K_a$ be a triangle with its three edges but no face.
In the edge order $(e_{01},e_{12},e_{02})$,
\[
 \partial_1=
 \begin{pmatrix}1&0&1\\1&1&0\\0&1&1\end{pmatrix},
 \qquad z=e_{01}+e_{12}+e_{02}.
\]
The matrix has rank two, so $Z_1=\langle z\rangle$.  Since $B_1=0$,
the quotient has two labels and exactly one nonzero exclusion label.
The uniform quotient-label entropy is one bit, and
$\Fill^{\mathrm{supp}}_a(z)=+\infty$.

At $b>a$, attach a face $\sigma$ with $\partial_2\sigma=z$.
Now $Z_1=B_1=\langle z\rangle$, the quotient is zero dimensional, and
\[
 \Fill^{\mathrm{supp}}_b(z)=1,\qquad d(z)=b,
 \qquad\mathsf H(L_1)=0.
\]
The uniform distribution on the \emph{single nonzero} label before
attachment would have entropy zero, not one bit; the one-bit statement
uses both quotient labels.
\end{example}

\Cref{fig:foundations-triangle} records the exact finite example above;
\cref{fig:foundations-costs} separates basis dependence from monotonicity
for a carried cycle.
\begin{figure}[tbp]
\centering
\begin{tikzpicture}
\begin{scope}[xshift=-3.5cm]
\draw[very thick,viewblue] (-1.4,0)--(1.4,0)--(0,1.9)--cycle;
\foreach \point in {(-1.4,0),(1.4,0),(0,1.9)} \fill[viewblue] \point circle (2pt);
\node at (0,-.45) {$K_a$: no face};
\node[vieworange] at (0,.65) {$z=e_{01}+e_{12}+e_{02}$};
\end{scope}
\begin{scope}[xshift=3.5cm]
\filldraw[fill=viewteal!15,draw=viewblue,very thick] (-1.4,0)--(1.4,0)--(0,1.9)--cycle;
\foreach \point in {(-1.4,0),(1.4,0),(0,1.9)} \fill[viewblue] \point circle (2pt);
\node at (0,.7) {$\sigma$};
\node at (0,-.45) {$K_b$: $\partial_2\sigma=z$};
\end{scope}
\draw[varrow] (-1.4,.85)--(1.4,.85) node[midway,above,font=\footnotesize]{attach one face};
\end{tikzpicture}
\medskip
\begin{tabular}{@{}lccccc@{}}
\toprule
Over $\F_2$ & $Z_1$ & $B_1$ & $\dim H_1$ & Labels / entropy & $\Fill^{\rm supp}(z)$\\\midrule
Before & $\langle z\rangle$ & $0$ & $1$ & $2$ / $1$ bit & $+\infty$\\
After & $\langle z\rangle$ & $\langle z\rangle$ & $0$ & $1$ / $0$ bits & $1$\\\bottomrule
\end{tabular}
\caption{One nonzero quotient label dies when one face is attached.
The same edge cycle is carried through the inclusion, and $d(z)=b$.
Entropy here uses all quotient labels, including zero.}
\label{fig:foundations-triangle}
\end{figure}
\begin{figure}[tbp]
\centering
\ViewHeading{(a) Support depends on the chosen basis}
\begin{tikzpicture}
\node[vbox,text width=5.5cm] (standard) at (-3.5,0)
  {Basis $((1,0),(0,1))$\\$(1,1)=1(1,0)+1(0,1)$\\support size $2$};
\node[vinput,text width=5.5cm] (changed) at (3.5,0)
  {Basis $((1,1),(0,1))$\\$(1,1)=1(1,1)+0(0,1)$\\support size $1$};
\end{tikzpicture}
\ViewHeading{(b) Filling of one carried cycle in compatible bases}
\begin{tikzpicture}[x=1.7cm,y=.7cm]
\draw[->,viewgray] (-.3,0)--(5.6,0) node[right] {$t$};
\draw[vieworange,very thick] (0,1)--(2.9,1);
\node[vieworange,above] at (1.3,1) {excluded: cost $+\infty$};
\draw[viewteal,very thick] (3,1)--(5,1);
\node[viewteal,above] at (4.2,1) {fillable};
\foreach \level/\name in {0/t_i,1/t_{i+1},2/t_{i+2},3/d(z),4/t_{j+1},5/t_{j+2}}
  {\fill[viewgray] (\level,0) circle (2pt);\node[below] at (\level,0) {$\name$};}
\draw[dashed,viewgray] (3,-.4)--(3,1.6);
\node[font=\footnotesize,align=center] at (2.5,-1.4)
 {Chosen start $t_i$ need not be a barcode birth.\\
  Once finite, minimum support cannot increase along inclusions.};
\end{tikzpicture}
\caption{Two separate qualifications on support filling. Changing bases
can change a finite support count. Monotonicity along a filtration uses
compatible bases and inclusions; it does not apply automatically to
arithmetic power maps or rescaled norms. Panel (b) is a schematic finite
threshold, not an additional computed example.}
\label{fig:foundations-costs}
\end{figure}

\section{Relative extension and its seed}
\label{sec:relative}

Let $(K,A)$ be a finite regular cellular pair with a finite-dimensional
cellular sheaf $\mathcal F$ over a field $\kk$, or a finite CW pair with
a finite-rank local system $\mathcal F$.  In either setting, use relative
cellular cochains as the kernel of restriction.  The degreewise short exact
sequence
\[
 0\to C^\bullet(K,A;\mathcal F)\to C^\bullet(K;\mathcal F)
       \to C^\bullet(A;\mathcal F)\to0
\]
gives the segment
\begin{equation}\label{eq:relative-sequence}
 H^0(K;\mathcal F)\xrightarrow{\mathrm{res}}H^0(A;\mathcal F)
               \xrightarrow{\delta}H^1(K,A;\mathcal F).
\end{equation}
Here an admissible local section is an element of $H^0(A;\mathcal F)$,
not an arbitrary assignment of values that may already violate the
relations on $A$.  For a network sheaf with specified coding data, this
criterion can describe extension of flows \cite{ghrist-hiraoka}; no
network interpretation is presumed for an arbitrary sheaf.

\begin{theorem}[Relative nonextension criterion]\label{thm:relative}
For $s\in H^0(A;\mathcal F)$,
\[
 s\notin\im(\mathrm{res})\iff\delta(s)\ne0.
\]
\end{theorem}

\begin{proof}
Exactness of \cref{eq:relative-sequence} says
$\ker\delta=\im(\mathrm{res})$.  Taking complements proves the claim.
\end{proof}

\subsection{A cochain presentation is a choice}

Choose a seed subspace $S\subseteq H^0(A;\mathcal F)$ and a linear
lifting $E:S\to C^0(K;\mathcal F)$ of its inclusion into
$C^0(A;\mathcal F)$.  Such a lift exists because degree-zero restriction
is surjective.  Since $s$ is a cocycle on $A$,
\[
 \widetilde\delta(s):=d^0E(s)\in Z^1(K,A;\mathcal F).
\]
The relative class $[\widetilde\delta(s)]$ is $\delta(s)$.
If $E'$ is another lift, then $E'-E$ takes values in relative degree-zero
cochains, and
\[
 d^0E'(s)-d^0E(s)\in B^1(K,A;\mathcal F).
\]
Thus the quotient seed is independent of the lift, although its displayed
cochain representative need not be.

This is the precise input supplied to the finite-response companion:
$\mathsf Z=Z^1(K,A;\mathcal F)$,
$\mathsf B=B^1(K,A;\mathcal F)$, $S$, and
$\widetilde\delta:S\to\mathsf Z$.  Endomorphisms preserving
$\mathsf B$ are additional marks; exactness does not select them.

\subsection{Naturality requires transport of the data}

Suppose a map of pairs with specified coefficient identifications induces
a morphism of the cochain short exact sequences.  Write $F_A$ for its
map on local sections and $F_{\mathrm{rel}}$ for its map on relative
cohomology.  Naturality gives
\[
 F_{\mathrm{rel}}\delta=\delta'F_A.
\]
A nonzero obstruction is preserved as nonzero when
$F_{\mathrm{rel}}$ is injective on that class; arbitrary maps can kill
classes.  A compatible pair homotopy equivalence for local systems gives
an isomorphism \cite[Section 3.H]{hatcher}.  A bare homotopy equivalence
without coefficient and seed identifications is not such a comparison.

\begin{example}[An additive cycle equation]\label{ex:additive}
For prescribed increments $a_j\in\F_Q$ around an oriented cycle, the
potential equations $u_{j+1}-u_j=a_j$ have a solution exactly when
$\sum_j a_j=0$.  Necessity follows by summing; sufficiency follows by
choosing $u_0$ and successively solving the equations.  This is an explicit
global consistency obstruction.  Identifying it with a particular
relative-sheaf seed would additionally require specifying that sheaf and
its local data.
\end{example}

\Cref{fig:foundations-relative} summarizes relative exactness, the choice
of lift, and the separate compatibility conditions for transport.
\begin{figure}[tbp]
\centering
\ViewHeading{(a) Extension is an exactness question}
\begin{tikzpicture}
\node[vbox] (global) at (-4.5,0) {$H^0(K;\mathcal F)$};
\node[vbox] (local) at (0,0) {$H^0(A;\mathcal F)$};
\node[vresult] (relative) at (4.5,0) {$H^1(K,A;\mathcal F)$};
\draw[varrow] (global)--(local) node[midway,above] {$\mathrm{res}$};
\draw[varrow] (local)--(relative) node[midway,above] {$\delta$};
\node[align=center] at (0,-.9)
 {$s$ extends $\Longleftrightarrow s\in\operatorname{im}(\mathrm{res})
   =\ker\delta\Longleftrightarrow\delta(s)=0$};
\end{tikzpicture}
\medskip
\begin{minipage}[t]{.47\linewidth}
\centering\ViewHeading{(b) Lifts need not be cocycles}
\begin{tikzpicture}
\node[vbox] (seed) at (0,0) {$s\in S$};
\node[vbox] (first) at (-1.55,-1.2) {$E(s)$};
\node[vbox] (second) at (1.55,-1.2) {$E'(s)$};
\node[vresult,text width=5.2cm] (quotient) at (0,-2.8)
  {$[d^0E(s)]=[d^0E'(s)]=\delta(s)$};
\draw[varrow] (seed)--(first);
\draw[varrow] (seed)--(second);
\draw[varrow] (first)--(quotient);
\draw[varrow] (second)--(quotient);
\node[font=\footnotesize] at (0,-3.6) {$d^0E'(s)-d^0E(s)\in B^1(K,A;\mathcal F)$};
\end{tikzpicture}
\end{minipage}\hfill
\begin{minipage}[t]{.50\linewidth}
\centering\ViewHeading{(c) Transport requires compatibility}
\begin{tikzpicture}
\node (source) at (0,0) {$H^0(A;\mathcal F)$};
\node (target) at (3.8,0) {$H^1(K,A;\mathcal F)$};
\node (sourcenew) at (0,-1.7) {$H^0(A';\mathcal F')$};
\node (targetnew) at (3.8,-1.7) {$H^1(K',A';\mathcal F')$};
\draw[varrow] (source)--(target) node[midway,above] {$\delta$};
\draw[varrow] (sourcenew)--(targetnew) node[midway,below] {$\delta'$};
\draw[varrow] (source)--(sourcenew) node[midway,left] {$F_A$};
\draw[varrow] (target)--(targetnew) node[midway,right] {$F_{\rm rel}$};
\node[align=center,text width=6cm,font=\footnotesize] at (1.9,-3)
 {A nonzero class stays nonzero only if the comparison does not kill it;
  injectivity on that class suffices.};
\end{tikzpicture}
\end{minipage}
\caption{Local sections must already satisfy the relations on $A$.
A lift to a degree-zero cochain is not automatically a global section.
Changing the lift changes the representative by a relative boundary;
transport between different data separately requires a compatible
cochain map and coefficient identifications.}
\label{fig:foundations-relative}
\end{figure}

\section{What does not follow from the qualitative tests}
\label{sec:limits}

\subsection{Fixed closed boundary spaces versus varying operators}

\begin{proposition}[Fixed finite-dimensional boundary spaces are closed]
\label{prop:closed}
Let $D:V\to W$ be a linear map between finite-dimensional real or complex
normed spaces.  If $z\notin\im D$, then
$\dist(z,\im D)>0$, so no sequence of values $Dc$ converges to $z$.
\end{proposition}

\begin{proof}
The image is a finite-dimensional subspace of $W$, hence closed.  Its
complement is open, so a ball of positive radius about $z$ misses it.
\end{proof}

This does \emph{not} say that quantitative degeneration requires
infinite-dimensional chain spaces.  For $D_t:\C\to\C$ given by
$D_t(c)=tc$, the minimum norm solving $D_t(c)=1$ is $1/|t|$ when
$t\ne0$ and $+\infty$ when $t=0$.  Every individual image is closed,
yet the minimum norm diverges as the operator varies toward zero.
The arithmetic companion similarly varies finite-dimensional boundary
operators with $x$ and $q$.  Its infinite index family is not an
infinite-dimensional completion of any one chain space.

\subsection{A checklist for a proposed bridge}

\Cref{fig:foundations-degeneration} contrasts the fixed-image obstruction
with the varying-operator example.
\begin{figure}[tbp]
\centering
\begin{minipage}[t]{.46\linewidth}
\centering\ViewHeading{(a) One fixed operator}
\begin{tikzpicture}
\draw[viewblue,very thick] (-2.5,0)--(2.5,0) node[below left] {$\operatorname{im}D$};
\node[vseed,label=above:$z$] (outside) at (.4,2) {};
\draw[dashed,vieworange] (.4,2)--(.4,0);
\draw[viewgray] (.4,.2)--(.6,.2)--(.6,0);
\node[align=center,font=\footnotesize] at (0,-1.2)
 {$z\notin\operatorname{im}D$\\$\operatorname{dist}(z,\operatorname{im}D)>0$};
\end{tikzpicture}
\end{minipage}\hfill
\begin{minipage}[t]{.50\linewidth}
\centering\ViewHeading{(b) A varying finite-dimensional operator}
\begin{tikzpicture}
\begin{axis}[viewaxis,width=6.7cm,height=4.5cm,xmin=0,xmax=2,ymin=0,ymax=9,
 xlabel={$t>0$},ylabel={minimum norm},xtick={0,1,2},ytick={1,4,8}]
\addplot[vieworange,domain=.12:2,samples=90] {1/x};
\node[font=\footnotesize,anchor=north east,align=right] at (rel axis cs:.95,.92)
 {$D_t(c)=tc$\\$\min_{tc=1}|c|=1/t$};
\end{axis}
\end{tikzpicture}
\end{minipage}
\caption{Closedness of each fixed image does not prevent minimum filling
norms from diverging as the operator varies. Panel (b) plots the exact
formula on $t>0$; at $t=0$ the equation has no solution and the cost is
$+\infty$. No infinite-dimensional completion is involved.}
\label{fig:foundations-degeneration}
\end{figure}

A transfer from a new geometric or filtered model to an arithmetic
spectrum must specify, and where necessary prove, the following:
\begin{enumerate}
 \item its source objects, coefficient systems, selected cycles or local
 sections, and the relevant comparison maps;
 \item its quantitative cost, including norms or bases, and the metric
 on the parameter or data space;
 \item a character assignment and its marked coordinates, with the
 approximation norm, constant, exponent, and infinitely-often condition;
 \item the precise relationship between source costs or classes and
 target data, including the direction of every implication;
 \item any metric estimate used to transfer dimension and the theorem
 giving the target dimension;
 \item an independently specified embedding and reach hypothesis if
 geometric offsets are claimed.
\end{enumerate}
These are admission requirements, not a theorem constructing the inputs.
A bi-Lipschitz bijection transfers Hausdorff dimension, but merely naming
one does not prove that it exists.

In the arithmetic example, a nonzero holonomy row makes a marked chain
fillable and a marked vertex section nonextendable.  At a trivial powered
character those conclusions reverse.  Therefore an arrow from ``chain
non-fillability'' to ``relative nonextension'' is not justified even by
their occurrence in the same explicit family.

\subsection{A formal conditional interface}

The following interface packages supplied data.  It neither constructs a
character assignment nor infers a norm from a qualitative obstruction.
Give $\T^d=(\R/\Z)^d$ its flat quotient metric
$d_{\T}(x,y)=(\sum_j\dist(x_j-y_j,\Z)^2)^{1/2}$.

\begin{definition}[Admitted dimension-and-nonextension data]
\label{def:admitted-bridge}
Fix an integer $d\geq1$ and a nonempty subset
$Y\subseteq\T^d\setminus\{0\}$ with the inherited metric.  An admitted
datum consists of:
\begin{enumerate}
 \item a metric space $(E,d_E)$ and a specified bijection $\Phi:E\to Y$
 with proved bounds
 \[
  a\,d_E(e,f)\leq d_{\T}(\Phi(e),\Phi(f))\leq b\,d_E(e,f)
          \quad(e,f\in E),\qquad 0<a\leq b<\infty;
 \]
 \item a connected finite CW complex $K$, a vertex $a_K$, and an
 epimorphism $\vartheta:\pi_1(K,a_K)\twoheadrightarrow\Z^d$;
 \item for every $e\in E$, a rank-one complex local system $\mathcal F_e$
 on $K$, a state $s_e\in H^0(\{a_K\};\mathcal F_e)$, and an isomorphism
 \[
  \psi_e:\mathcal F_e\xrightarrow{\ \cong\ }
                  \mathcal L_{\Phi(e)}^\vartheta,
  \qquad (\psi_e)_{a_K}(s_e)=1.
 \]
\end{enumerate}
Here $\mathcal L_y^\vartheta$ has the marked fibre $\C$ and holonomy
$g\mapsto\exp(2\pi\mathrm i\langle y,\vartheta(g)\rangle)$.
The subset $Y$, including any approximation norm, exponent, and tolerance
used to define it, is part of the data.  No continuity in $e$ of the
coefficient identifications is assumed or needed for the conclusions below.
\end{definition}

\begin{proposition}[Conditional dimension and nonextension transport]
\label{prop:admitted-bridge}
For data in \cref{def:admitted-bridge},
\[
 \dim_{\mathrm H}(E,d_E)=\dim_{\mathrm H}(Y,d_{\T}),\qquad
 \delta_e(s_e)\ne0\quad\text{in }H^1(K,\{a_K\};\mathcal F_e)
                         \quad(e\in E).
\]
In particular, every selected state fails to extend globally.  If an
independent theorem gives $\dim_{\mathrm H}Y=h$, then
$\dim_{\mathrm H}E=h$.
\end{proposition}

\begin{proof}
The upper metric bound makes $\Phi$ Lipschitz, so covers of $E$ push
forward with diameters multiplied by at most $b$.  The lower bound makes
$\Phi^{-1}$ Lipschitz with constant $a^{-1}$.  Applying the definition
of Hausdorff measure in both directions gives equality of dimensions.

Separately, let $y=\Phi(e)\ne0$.  Some coordinate of $y$ is nonzero
modulo $\Z$, and surjectivity of $\vartheta$ supplies a loop with
nonidentity scalar holonomy.  A global section of the rank-one system
is fixed by every holonomy, so it must vanish.  The relative exact
sequence therefore injects the marked vertex fibre $\C$ into
$H^1(K,\{a_K\};\mathcal L_y^\vartheta)$, sending $1$ to a nonzero
class.  Naturality with respect to $\psi_e$ identifies this class with
$\delta_e(s_e)$.  Finally apply \cref{thm:relative}.
\end{proof}

The two conclusions use different inputs: metric estimates give dimension,
whereas coefficient and state identifications give nonextension.  Neither
proof gives a quantitative filling cost, an offset radius, or endomorphism
marks.  If $Y$ is specified by a cube-based periodic locus, first take its
image in $\T^d$ and then remove the trivial character.  Removing only one
corner of $[0,1]^d$ before taking the quotient does not enforce $Y\cap\{0\}
=\varnothing$.  The arithmetic companion's normalization appendix treats
the earlier sup-norm locus and the chord-norm locus as distinct admissible
targets, not as equal sets.

\subsection{Dimension, codimension, and geometric radius}

The two conclusions of the conditional interface use different inputs,
as shown in \cref{fig:foundations-bridge}.
\begin{figure}[tbp]
\centering
\begin{tikzpicture}
\node[vinput,text width=5.6cm] (metric) at (-3.55,0)
  {\textbf{Metric input}\\A specified bijection $\Phi:E\to Y$\\
   $a\,d_E\leq d_{\T}\circ(\Phi,\Phi)\leq b\,d_E$};
\node[vinput,text width=5.6cm] (coeff) at (3.55,0)
  {\textbf{Coefficient and state input}\\
   $\psi_e:\mathcal F_e\cong\mathcal L_{\Phi(e)}^\vartheta$\\
   $(\psi_e)_{a_K}(s_e)=1$, $\Phi(e)\neq0$};
\node[vbox,text width=5.6cm] (dimproof) at (-3.55,-1.95)
  {Lipschitz cover comparisons\\in both directions};
\node[vbox,text width=5.6cm] (obproof) at (3.55,-1.95)
  {Surjective $\vartheta:\pi_1(K)\twoheadrightarrow\Z^d$\\
   Nontrivial rank-one holonomy; exactness};
\node[vresult,text width=5.6cm] (dim) at (-3.55,-3.6)
  {$\dim_{\rm H}E=\dim_{\rm H}Y$};
\node[vresult,text width=5.6cm] (ob) at (3.55,-3.6)
  {$\delta_e(s_e)\neq0$};
\draw[varrow] (metric)--(dimproof);
\draw[varrow] (dimproof)--(dim);
\draw[varrow] (coeff)--(obproof);
\draw[varrow] (obproof)--(ob);
\node[font=\footnotesize,align=center] at (0,-4.65)
 {Independent conclusions from different supplied inputs.\\
  Neither route constructs a filling norm, an offset radius, or response marks.};
\end{tikzpicture}
\caption{The conditional interface has two proof routes, not a deduction
of dimension from nonextension or conversely. An independent theorem is
still needed to evaluate $\dim_{\rm H}Y$.}
\label{fig:foundations-bridge}
\end{figure}

The general spectrum in the arithmetic companion has dimension
$d-\eta$ after choosing $\alpha=(1+\eta)/(d-\eta)$, where
$0<\eta<d$.  Its $d=2$, $\eta=1/(64\pi^2)$ example retains the earlier
decimal $1.9984$ as a rounded Hausdorff dimension.  The codimension is
$\eta$, not $d-\eta$.  Both are separate from the product radii and
offset radius.

No claim of a universal fractional manifold dimension, geometric buffer
width, physical information law, or independently derived numerical
constant is retained.  A second occurrence of the same scalar in another
model would not establish a common mechanism without an explicit map and
comparison theorem.  The detailed Diophantine and reach calculations now
belong to the core companion rather than being repeated here.

\section{Conclusion and finite verification}

There are three precise conclusions: boundary writing produces exactly
the boundary subspace; finite inclusion filtrations make that solvability
test scale sensitive; and relative exactness detects failure of a local
section to extend.  Only the first, with a supplied uniform distribution,
gives the finite-label entropy stated here.

The quantitative companion constructs one additional model in which
normed filling costs and arithmetic recurrence admit a complete
comparison.  It does not close the general bridge problem for arbitrary
filtered complexes.  The response-module companion supplies a different
finite operation once its quotient, seed, and marks are fixed.

The script \path{scripts/verify_restructured_series.py} checks the
triangle boundary relation and ranks over $\F_2$, the associated label
counts, and the displayed finite-response examples.  These calculations
are finite consistency checks, not substitutes for the proofs.

\appendix
\section{General geometric transport below reach}
\label{sec:general-reach}

Let $M\subseteq\R^m$ be nonempty and closed.  Its reach is the supremum
of all $\rho>0$ such that every point at distance less than $\rho$ from
$M$ has a unique nearest point.  Write
$M^r=\{y:\dist(y,M)\leq r\}$.  The following standard consequence of
positive reach is included with a proof to distinguish its generality
from the product-torus calculation; see Federer \cite[Section 4]{federer}
for the underlying geometric framework.

\begin{lemma}[General sub-reach strong deformation retraction]
\label{lem:general-reach}
Suppose $\reach(M)>0$.  On
$V=\{y:\dist(y,M)<\reach(M)\}$, nearest-point projection $p:V\to M$
is continuous.  For every finite $0\leq r<\reach(M)$, its restriction
$p_r:M^r\to M$ gives a strong deformation retraction through
\[
 H_r(y,t)=(1-t)y+t p(y),\qquad 0\leq t\leq1.
\]
For $0\leq r\leq R<\reach(M)$, with $R$ finite, inclusion
$j_{r,R}:M^r\hookrightarrow M^R$ satisfies $p_Rj_{r,R}=p_r$.
Every $A\subseteq M$ is fixed pointwise by these homotopies, so they
are also homotopies of pairs relative to $A$.
\end{lemma}

\begin{proof}
A nearest point exists for every $y\in\R^m$: a minimizing sequence in
the nonempty closed set $M$ is bounded and has a convergent subsequence
in $M$.  Uniqueness on $V$ follows from the reach hypothesis.
To prove continuity, let $y_n\to y$ in $V$.  The points $p(y_n)$ are
bounded, since
\[
 \norm{p(y_n)-y}\leq\dist(y_n,M)+\norm{y_n-y}.
\]
For every convergent subsequence with limit $z$, closedness puts $z$ in
$M$, and continuity of distance gives
$\norm{z-y}=\dist(y,M)$.  Uniqueness forces $z=p(y)$.  If the full
sequence did not converge to $p(y)$, a bounded subsequence staying a
positive distance away would have a convergent subsequence with another
limit, a contradiction.

Set $z=p(y)$ and $D=\norm{y-z}$.  For $w=(1-t)y+tz$ and any $u\in M$,
\[
 \norm{w-u}\geq\norm{y-u}-\norm{y-w}\geq D-tD=(1-t)D,
 \qquad \norm{w-z}=(1-t)D.
\]
Hence $\dist(w,M)=(1-t)D\leq r$, and $p(w)=z$ by uniqueness.
Continuity of $p$ makes $H_r$ continuous; it starts at the identity,
ends at $p_r$, and fixes $M$ at every time.  The same nearest-point
function defines every $p_r$, proving compatibility under inclusion and
the assertion about subsets $A$.
\end{proof}

This lemma does not require $M$ to be a manifold, a product, or compact.
If a local system and a marked state on $M$ are independently supplied,
the compatible pullbacks along $p_r$ provide the setting for relative
cohomology transport by pair homotopy invariance.  The lemma itself does
not choose those coefficients or states.  Nor does it provide a specified
finite cellularization, a compatible sampling filtration, a Diophantine
parameterization, or a normed filling model.  Those additional structures
must be stated separately; the arithmetic companion supplies them only
for its explicit carrier construction.

\section{Visual dictionary for the three manuscripts}
\label{sec:visual-dictionary}

\Cref{fig:series-dictionary} is a reference guide to the inputs and indices
used in the three papers. Its rows describe different constructions;
they do not assert comparison maps between them.
\begin{figure}[htbp]
\centering
\ViewHeading{(a) Inputs, operations, and outputs}
\begin{tabular}{@{}>{\raggedright\arraybackslash}p{.27\linewidth}
 >{\raggedright\arraybackslash}p{.30\linewidth}
 >{\raggedright\arraybackslash}p{.34\linewidth}@{}}
\toprule
Supplied input & Operation & Output\\\midrule
Finite $\F_Q$ chain complex & Quotient by boundaries & Labels in $H_n$; entropy only after choosing a distribution\\[4pt]
Marked normed arithmetic family & Minimize filling norm; select recurrence & $\Fill_{q,x}(1)$ and $R_{\alpha,C}$\\[4pt]
Finite quotient, seed, marks & Close under operator words and linear span & Minimal invariant response $\Nex$\\
\bottomrule
\end{tabular}
\medskip
\ViewHeading{(b) Four different indices}
\begin{tabular}{@{}cll@{}}
\toprule
Index & What changes & Comparison direction\\\midrule
$t$ & Complex in an inclusion filtration & $K_{t_i}\hookrightarrow K_{t_j}$, $t_i\leq t_j$\\[3pt]
$q$ & Powered character $\mathcal L_{qx}$ & $\mathcal C_x(kq)\longrightarrow\mathcal C_x(q)$\\[3pt]
$r$ & Offset; character stays fixed & $H^1(M^R,\{a\})\longrightarrow H^1(M^r,\{a\})$\\[3pt]
$m$ & Subspace inside one fixed quotient & $N^{(m)}\subseteq N^{(m+1)}$\\
\bottomrule
\end{tabular}
\caption{A reading dictionary for the series. Coefficients in the offset
row are the compatible pulled-back local systems and are suppressed only
to shorten the table. There is no identification of $t,q,r$, or the
closure index $m$.}
\label{fig:series-dictionary}
\end{figure}

\bibliographystyle{plain}
\bibliography{series_references}
\end{document}